\chapter{RESULTADOS ESPERADOS}

% Nesta seção deve ser descrito o que se espera de resultado final após a conclusão do trabalho.
Espera-se, ao final deste trabalho, o desenvolvimento de uma aplicação capaz de realizar testes comparativos entre diferentes implementações de \glspl{bvh} em um ambiente de renderização híbrida utilizando técnicas básicas de \textit{ray tracing}. A aplicação deverá permitir a execução controlada de experimentos, coleta de métricas de desempenho e análise do comportamento das estruturas em diferentes tipos de cena e condições de \textit{hardware}.

Adicionalmente, identificar características, vantagens e limitações de cada algoritmo de \gls{bvh} analisado. A partir dos dados obtidos, pretende-se determinar quais estruturas apresentam melhor desempenho em cenários específicos, considerando fatores como quantidade de objetos dinâmicos, complexidade geométrica da cena, custo de atualização, consumo de memória e utilização da \gls{gpu} e da \gls{cpu}.

Além disso, espera-se compreender como diferentes abordagens de \glspl{bvh} se comportam em aplicações de renderização em tempo real, permitindo identificar quais algoritmos são mais adequados para diferentes contextos computacionais.

Por fim, espera-se que os resultados obtidos contribuam para o entendimento do impacto das \glspl{bvh} no desempenho de aplicações que utilizam \textit{ray tracing}, servindo como base para estudos futuros relacionados à otimização de estruturas espaciais e técnicas de aceleração gráfica.